% MAS1004 Assignment 2 long report, English edition prepared 10 November 2026.
% Compile with XeLaTeX. The course style is supplied unchanged.
\documentclass[11pt]{article}
\usepackage{mas1004}
\IfFontExistsTF{Noto Serif CJK KR}{}{%
  \IfFontExistsTF{AppleMyungjo}{%
    \setmainhangulfont[AutoFakeBold=2]{AppleMyungjo}
    \setsanshangulfont{Apple SD Gothic Neo}
    \setmonohangulfont{AppleMyungjo}
  }{}%
}
\hypersetup{pdftitle={Assignment 2 Long report},pdfauthor={김서영}}
\reporttitle{Assignment 2 Long report}
\studentname{김서영}
\studentid{20251120}

\begin{document}
\makeheader
{\small Leaderboard nickname \texttt{SoSleepy}\hfill
Selected releases \texttt{track-a-2} and \texttt{track-b-2}}\par

\section{Scope and experimental purpose}

This assignment classifies photographs as Seoul, Busan, Tokyo, Osaka, or Taipei.
We first measured how much the validation score changed when the starter recipe
was trained again with a different seed. After completing the baseline with
seeds 0 and 1, the student reviewed mistakes from all five cities and selected
56 training photographs to exclude. With the same two seeds, the mean
validation score was 50.14\% with all training photographs and 49.26\% after
excluding the selected photographs. The student then proposed restoring all
training photographs and masking small random regions during training.
The mean score for that setting was 49.30\%.
The student later selected balanced sampling. A paired comparison on the same
Mac and seed 0 scored 49.51\% with the starter sampling and 50.91\% with
balanced sampling, an observed increase of 1.40 percentage points.

The report covers eight full Track A training runs across the recorded CUDA
and MPS conditions and three full Track B training runs. The first Track B model scored 48.95\% in its
training-time validation. A subsequent paired comparison used the same MPS
device and seed 0, with and without seven photographs supplied by the student.
The scores were 48.61\% and 48.74\%, an observed increase of 0.13 percentage
points. Each of these two conditions was run once.

On 10 November, the student selected reuse of the completed balanced Track A
model for Track B. Its full ONNX/CPU validation score was 50.91\%, and it
was published as \texttt{track-b-2}. The leaderboard updated at 17:45 KST
reported 42.03\% and rank 14 among 27 scored models. Reusing this model
added no training run. Student-selected MobileNetV3-Large experiments were
still in progress when this revision was prepared, so their results are
reserved for a subsequent revision.

After training, we performed an optional occlusion analysis on ten validation
photographs using the first Track B model, \texttt{track-b-1}. The student\textquotesingle{}s question about an
unexpected prediction change led to an additional comparison of replacement
colors and blur at the same image coordinates.

\begin{center}\small
\begin{tabular}{@{}p{0.25\linewidth}p{0.70\linewidth}@{}}
\toprule
Work & Completed work or recorded status \\
\midrule
Track A baseline & Full training with seeds 0 and 1, followed by model format checks \\
Photo review & Mistakes from all five cities reviewed, with 56 exclusions selected from 480 training candidates \\
Track A changes & Exclusion and Random Erasing with two CUDA seeds, followed by balanced sampling and a matched MPS control with seed 0 \\
Track A release & \texttt{track-a-2}, balanced sampling using the full training selection, selected and checked \\
Track B experiments & Three full training runs, including a paired comparison of seven additional photographs \\
Track B release & Balanced model reused as \texttt{track-b-2}, with full ONNX/CPU validation of 50.91\% \\
Leaderboard & Track B scored 42.03\%, rank 14 of 27, in the 17:45 KST update on 10 November. Track A had no matching entry \\
\bottomrule
\end{tabular}
\end{center}

\section{Data and Track A training conditions}

The baseline used all starter training photographs. The following training
counts come from the baseline \texttt{result.json}, and the validation counts
are row totals of the saved confusion matrices. Counts for the exclusion
condition are reported with that experiment.

\begin{center}\small
\begin{tabular}{lrr}
\toprule
City & Training photographs & Validation photographs \\
\midrule
Seoul & 11,274 & 2,435 \\
Busan & 5,191 & 1,390 \\
Tokyo & 13,347 & 3,881 \\
Osaka & 7,200 & 2,128 \\
Taipei & 7,245 & 2,243 \\
\midrule
Total & 44,257 & 12,077 \\
\bottomrule
\end{tabular}
\end{center}

The model was ResNet18 initialized with ImageNet pretrained weights, with an
input size of 224 $\times$ 224. Batch size 64 and 3,906 training steps resulted
in 249,984 photograph presentations per run. All eight full Track A result files
recorded \texttt{official\_budget=true}.

The starter recipe applies random resized cropping, horizontal flipping, and
color jitter. The crop area ranges from 0.35 to 1.0 of the original image area.
Brightness, contrast, and saturation jitter are each 0.2. Training uses AdamW
with learning rate $3\times10^{-4}$ and weight decay 0.05, together with
OneCycleLR and label smoothing of 0.1. The recipe is preserved in
\path{outputs/recipe_snapshot.py}.

Both baseline runs used CUDA. The recipe and photo selection were identical,
and the seed changed from 0 to 1. Both recipe hashes begin with
\texttt{17ac485223fa}. The recorded commits were \texttt{a23f9e6} and
\texttt{0dd327b}, respectively.

The validation score is the mean of the five per-city accuracies. Each city
therefore contributes equally, despite differences in the number of
validation photographs.

\section{Baseline results}

\begin{center}\small
\begin{tabular}{llrr}
\toprule
Run tag & Seed & Training time (min) & Validation (\%) \\
\midrule
\texttt{baseline} & 0 & 37.3 & 50.44 \\
\texttt{baseline} & 1 & 36.3 & 49.84 \\
\midrule
Two-run mean & & & 50.14 \\
\bottomrule
\end{tabular}
\end{center}

The recorded training times were 2,238.2 and 2,179.5 seconds. Data preparation,
validation, and model export required additional time. The per-city scores
below come from these same validation runs.

\begin{center}\small
\begin{tabular}{lrr}
\toprule
City & Seed 0 accuracy (\%) & Seed 1 accuracy (\%) \\
\midrule
Seoul & 57.74 & 56.26 \\
Busan & 39.42 & 38.42 \\
Tokyo & 63.85 & 61.40 \\
Osaka & 36.04 & 38.91 \\
Taipei & 55.15 & 54.21 \\
\bottomrule
\end{tabular}
\end{center}

The absolute difference between the two overall scores was 0.60 percentage
points. This illustrates variation from random choices under the same training
method. Small improvements should be compared against a baseline using the
same seed and checked with repeated runs when needed. Two runs provide
limited information about the full distribution of training variability or
statistical significance.

Tokyo had the highest accuracy in both runs. Busan and Osaka had lower
accuracies than the other cities. The agent first presented mistake sheets
for Busan and Osaka, and the student inspected them and recorded observations.

The Colab check outputs for both ONNX models reported 44.7 MB, 11.17 million
parameters, and 1.81 billion multiply-adds for a 224 $\times$ 224 input.
Both printed \texttt{format and limits: OK}. The original outputs are included
in the appendix.

\section{Photo observations and hypotheses}

\subsection{How the mistake sheets were reviewed}

The command \texttt{curate.py mistakes} applied the seed 0 ONNX model to the
validation photographs and selected 48 high-confidence mistakes per city.
The five sheets contain 240 photographs in total. The city in a sheet\textquotesingle{}s name
is the ground-truth label. The city after the arrow under each photograph is
the prediction, followed by its softmax output probability.

These sheets focus on high-confidence mistakes. We also examined the complete
validation confusion matrices saved by \texttt{official.py}. The sheets used
CPU inference in ONNX Runtime, while the validation score tables use the
values recorded by the official training code.

\subsection{Similar urban scenes in Busan and Seoul}

The student noticed streetscapes that appeared common to Busan and Seoul.
During the review on 7 November, the student explained this in the following
statement, translated from Korean.

\begin{quote}\small
Without the sea or a maritime feel, I do not think it is easy to tell whether
a photograph is from Busan or Seoul.
\end{quote}

An earlier observation pointed to cars parked along a residential street in
Busan photograph 8. The student hypothesized that ordinary roads and buildings
shared by the two cities could contribute to confusion.

Of the 48 mistakes on the Busan sheet, 34 were predicted as Seoul. The full
validation confusion matrices showed the following pattern.

\begin{center}\small
\begin{tabular}{lrrrr}
\toprule
Seed & \shortstack{Busan\\photographs} & \shortstack{Busan\\errors} &
\shortstack{Predicted\\Seoul} & \shortstack{Seoul share of\\Busan errors (\%)} \\
\midrule
0 & 1,390 & 842 & 442 & 52.49 \\
1 & 1,390 & 856 & 463 & 54.09 \\
\bottomrule
\end{tabular}
\end{center}

More than half of the Busan errors in each run were predictions of Seoul.
The causal role of visual similarity remains a hypothesis for further case
review and measurement. The student\textquotesingle{}s initial impression about street parking
in Japanese residential areas also requires a broader set of examples.

\subsection{Utility poles and overhead wires in Busan photograph 48}

The student noticed utility poles and wires in Busan photograph 48 and thought
their association with Japan might be related to the mistake. Its ground-truth
label is Busan, and the model predicted Tokyo with output probability 0.85.

Overhead wires also exist in Korea. The Seoul Metropolitan Government\textquotesingle{}s
published undergrounding plan describes sections where electricity and
telecommunications lines and poles are to be moved underground. The presence
of wires should be assessed alongside other visual cues, and their actual
influence on this prediction requires measurement. The source is listed in
the references.

\subsection{Similar scenes in Osaka and Tokyo}

The student\textquotesingle{}s initial observations on the Osaka sheet included shrines,
houses, and residential streets with no visible sidewalks. On 8 November,
the student explained that Osaka and Tokyo might be confused because both
are Japanese cities with similar residential and shrine scenes.

Among Osaka errors in the baseline confusion matrices, Tokyo predictions
accounted for 855/1,361, or approximately 62.82\%, with seed 0, and 783/1,300,
or approximately 60.23\%, with seed 1. The relationship between the shared
scenes and these mistakes was recorded as a photo-based hypothesis.

\subsection{Food photographs in Seoul and vehicle photographs in Tokyo}

The student found it difficult to identify the city from Seoul photographs
4 and 10, which mainly showed food. Subway and road scenes also appeared to
share features across cities. For Seoul photograph 24, the student suggested
that reflective paving blocks and an orderly street might evoke Japan.
National differences in paving materials and the model\textquotesingle{}s actual use of these
features remained questions for verification.

For Tokyo photographs 4 and 8, the student similarly thought that a vehicle
close-up could contain few city-specific cues. Subways and busy commercial
streets could also resemble Osaka. The labels and predictions were checked
against the sheet CSV files.

\begin{center}\small
\begin{tabular}{lllr}
\toprule
Sheet and number & Ground truth & Prediction & Output probability \\
\midrule
Seoul 4 & Seoul & Taipei & 0.96 \\
Seoul 10 & Seoul & Osaka & 0.91 \\
Seoul 24 & Seoul & Tokyo & 0.88 \\
Tokyo 4 & Tokyo & Seoul & 0.96 \\
Tokyo 8 & Tokyo & Osaka & 0.95 \\
\bottomrule
\end{tabular}
\end{center}

Seoul photographs 11 and 46 were checked again after their labels and
prediction directions were initially misread. Both have Seoul labels, with
Tokyo and Taipei predictions respectively. The interpretation was corrected
by distinguishing people\textquotesingle{}s positions on an escalator from the side of the
road used by vehicles.

\subsection{Cosplay, commercial spaces, and waterfront scenes in Taipei}

The student explained that Taipei\textquotesingle{}s appearance and culture were unfamiliar.
The cosplay in photograph 41 looked like a Japanese manga character and might
evoke Tokyo. The crowded commercial spaces in photographs 1 and 33 reminded
the student of Hongdae and Seongsu in Seoul. Water and waterfront structures
in photograph 10 suggested a sea and breakwater, which evoked Busan. Simple
nature photographs also seemed to provide few clues for identifying a city.

\begin{center}\small
\begin{tabular}{lllr}
\toprule
Sheet and number & Ground truth & Prediction & Output probability \\
\midrule
Taipei 1 & Taipei & Seoul & 0.96 \\
Taipei 10 & Taipei & Busan & 0.91 \\
Taipei 33 & Taipei & Seoul & 0.85 \\
Taipei 41 & Taipei & Tokyo & 0.83 \\
\bottomrule
\end{tabular}
\end{center}

The character\textquotesingle{}s identity, the type of waterfront, and the features used by
the model require further verification. These observations motivated a review
of photographs with limited city background. The student\textquotesingle{}s original Korean
observations and the review process are preserved in the lab notebook at the
end of the report.

\section{Excluding 56 training photographs selected by the student}

\subsection{Selection process and experimental scope}

Following the observations about food in Seoul photographs and vehicles in
Tokyo photographs, the student first chose to curate training photographs
with food or vehicle close-ups. The agent applied the baseline seed 0 model
to every fourth training photograph, a sample of 11,065. Photographs receiving
a low probability for their ground-truth city were collected into 96
candidates per city, or 480 in total. Each city had sheets 00 and 01, with
photographs numbered from 1 to 48 on each sheet.

The student directly selected the numbers listed below. The actual selections
also included birds, signs, sky, book covers, and flowers. We therefore
recorded the experiment as curation of the photographs individually selected
from the candidate sheets. The sheet names, numbers, and paths were retained
without assigning an inferred reason to each photograph.

\begin{center}\small
\begin{tabular}{llp{0.56\linewidth}r}
\toprule
City & Sheet & Numbers selected by the student & Count \\
\midrule
Seoul & 00 & 5, 10, 14, 45 & 4 \\
Seoul & 01 & 2, 28, 30, 37, 39 & 5 \\
Busan & 00 & 1, 16, 31, 45 & 4 \\
Busan & 01 & 4, 24, 29, 38, 40 & 5 \\
Tokyo & 00 & 4, 11, 16, 19, 28, 34, 42 & 7 \\
Tokyo & 01 & 1, 3, 8, 10, 11, 15, 16, 21, 27, 33, 37, 42, 48 & 13 \\
Osaka & 00 & 18, 24, 28, 38, 39 & 5 \\
Osaka & 01 & 33 & 1 \\
Taipei & 00 & 2, 14, 30, 39, 41, 45, 46 & 7 \\
Taipei & 01 & 3, 11, 20, 23, 41 & 5 \\
\midrule
Total & & & 56 \\
\bottomrule
\end{tabular}
\end{center}

The numbers and paths were checked against the sheet CSV files to create
\texttt{curation\_plan.csv}. Exactly these 56 photographs, approximately
0.13\% of the training set, were excluded from the full training inventory
on Colab. The experiment measures the effect of this selection from the
low-confidence candidate sheets.

\begin{center}\small
\begin{tabular}{lrrr}
\toprule
City & Original training set & Excluded & Retained \\
\midrule
Seoul & 11,274 & 9 & 11,265 \\
Busan & 5,191 & 9 & 5,182 \\
Tokyo & 13,347 & 20 & 13,327 \\
Osaka & 7,200 & 6 & 7,194 \\
Taipei & 7,245 & 12 & 7,233 \\
\midrule
Total & 44,257 & 56 & 44,201 \\
\bottomrule
\end{tabular}
\end{center}

Both runs used the baseline training recipe, with hash prefix
\texttt{17ac485223fa}. The selection hash prefix was \texttt{9ef44cac68cb}.
The actual selection, exclusion record, original plan, and hash summary were
exported from Colab and stored with the result ZIP. Checks covered duplicate
paths and overlap between retained and excluded photographs. Combining both
lists reproduced the original training count and inventory hash.

\subsection{Results paired by seed}

\begin{center}\small
\begin{tabular}{lrrr}
\toprule
Run & All photographs (\%) & Exclude 56 (\%) & Change (pp) \\
\midrule
Seed 0 & 50.44 & 49.37 & -1.07 \\
Seed 1 & 49.84 & 49.14 & -0.70 \\
\midrule
Two-run mean & 50.14 & 49.26 & -0.89 \\
\bottomrule
\end{tabular}
\end{center}

Before rounding, the means were 0.5014 and 0.49255. Their difference was
-0.00885, or -0.885 percentage points. Values were calculated from the saved
result JSON files and then rounded to two decimal places. Subtracting the
displayed means can therefore differ in the last digit from the displayed
change.

\begin{center}\small
\begin{tabular}{lrrr}
\toprule
City & All-photo mean (\%) & Exclude-56 mean (\%) & Change (pp) \\
\midrule
Seoul & 57.00 & 57.66 & +0.66 \\
Busan & 38.92 & 36.59 & -2.34 \\
Tokyo & 62.63 & 64.49 & +1.86 \\
Osaka & 37.48 & 36.11 & -1.37 \\
Taipei & 54.68 & 51.43 & -3.26 \\
\bottomrule
\end{tabular}
\end{center}

Training took 2,211.8 seconds for seed 0 and 2,220.7 seconds for seed 1,
approximately 36.9 and 37.0 minutes. Both ONNX models passed the format and
limit checks. The run commits were \texttt{9cd86e1} and \texttt{a2ac080}.

The overall score decreased with both seeds. Mean accuracy increased for
Seoul and Tokyo and decreased for Busan, Osaka, and Taipei. The measured
conclusion concerns this particular set of 56 exclusions, these two seeds,
and the public validation set. Further experiments would be needed to
separate individual-photo effects or generalize to other curation rules.
Changing the selection length also changes the official sampler\textquotesingle{}s random
presentation order, including when the same seed is used.

\section{Random Erasing proposed by the student}

After reviewing both exclusion runs, the student made the following request,
translated from Korean.

\begin{quote}\small
What about using all the photographs and masking some part of each one?
\end{quote}

After clarification, the student chose to mask small random regions during
training so that the model could learn from multiple cues. The agent
implemented the following initial comparison using torchvision\textquotesingle{}s
\texttt{RandomErasing}.

\begin{center}\small
\begin{tabular}{@{}p{0.27\linewidth}p{0.68\linewidth}@{}}
\toprule
Item & Configuration \\
\midrule
Experiment tag & \texttt{erase15} \\
Training photographs & All 44,257 photographs restored in the original order \\
Application probability & 0.5 each time a training photograph is presented \\
Erased area fraction & 0.02 to 0.15 of the transformed image area \\
Rectangle aspect ratio & 0.3 to 3.3 \\
Fill RGB & (0.485, 0.456, 0.406), which becomes 0 after ImageNet normalization \\
Other transformations & Random resized crop with area 0.35 to 1.0, horizontal flip, and color jitter \\
Comparison & Paired comparisons against the all-photo baseline with seeds 0 and 1 \\
\bottomrule
\end{tabular}
\end{center}

The student selected the masking approach. The agent chose the probability,
area range, and other numerical settings for this first comparison. The
hypothesis was that learning from the remaining city cues after a small
region was hidden might improve validation performance. The comparison
restored the baseline\textquotesingle{}s complete training inventory and added the erasing
augmentation.

Erasing was applied after \texttt{ToTensor}, followed by normalization in the
official code. With the assignment\textquotesingle{}s PyTorch 2.6.0 and torchvision 0.21.0,
checks covered recipe rules, actual transformations, the official photo
loader, and forward and backward passes on synthetic data. The recipe
snapshot is stored in \path{outputs/erase15_recipe.py}. Full training and
validation were then completed on Colab with seeds 0 and 1.

\subsection{Measured effect of erasing}

\begin{center}\small
\begin{tabular}{lrrr}
\toprule
Run & Baseline (\%) & Random Erasing (\%) & Change (pp) \\
\midrule
Seed 0 & 50.44 & 49.16 & -1.28 \\
Seed 1 & 49.84 & 49.44 & -0.40 \\
\midrule
Two-run mean & 50.14 & 49.30 & -0.84 \\
\bottomrule
\end{tabular}
\end{center}

\begin{center}\small
\begin{tabular}{lrrr}
\toprule
City & Baseline mean (\%) & Erasing mean (\%) & Change (pp) \\
\midrule
Seoul & 57.00 & 57.09 & +0.09 \\
Busan & 38.92 & 37.77 & -1.15 \\
Tokyo & 62.63 & 63.22 & +0.59 \\
Osaka & 37.48 & 35.20 & -2.28 \\
Taipei & 54.68 & 53.24 & -1.45 \\
\bottomrule
\end{tabular}
\end{center}

Each seed used the same training count and inventory order as the baseline,
with matching batch size, step count, presentation budget, and validation
photographs. Both runs had recipe hash prefix \texttt{4c2c34345077} and
all-photo selection hash prefix \texttt{1a1ca6c18823}. The overall and per-city
scores were recalculated from the saved confusion matrices and checked
against the original results.

Training took 2,208.3 and 2,207.2 seconds, approximately 36.8 minutes each.
Both runs used commit \texttt{c6bf059}. They finished at 23:57 KST on
8 November and 00:37 KST on 9 November 2026. Both ONNX models passed the
format and limit checks.

Under this probability and area range, the overall score decreased with both
seeds. Mean accuracy increased for Seoul and Tokyo and decreased for Busan,
Osaka, and Taipei. The comparison measures the performance of this particular
configuration. Establishing which cues the model used, or the cause of the
decrease, requires additional analysis. Further student interpretations can
be recorded separately from these numerical comparisons.

\subsection{Comparison of the initial CUDA conditions}

\begin{center}\small
\begin{tabular}{p{0.44\linewidth}rrr}
\toprule
Condition & Seed 0 (\%) & Seed 1 (\%) & Mean (\%) \\
\midrule
Baseline recipe, all photographs & 50.44 & 49.84 & 50.14 \\
Student-selected 56-photo exclusion & 49.37 & 49.14 & 49.26 \\
All photographs with Random Erasing & 49.16 & 49.44 & 49.30 \\
\bottomrule
\end{tabular}
\end{center}

The baseline recipe had the highest two-seed mean among the three measured
conditions. The exclusion and erasing runs are retained as measurements of
the respective changes under their recorded settings.

\section{Balanced sampling with a matched local control}
\label{sec:balanced}

\subsection{Choice and implementation}

After seeing balanced sampling among other participants\textquotesingle{}
settings, the student explicitly selected it as the next experiment.
The agent explained the training imbalance and the scoring rule. Tokyo has
13,347 training photographs and Busan has 5,191, so their shares in the
starter selection are approximately 30.16\% and 11.73\%. The validation
metric gives each of the five cities the same weight.

The new recipe gives each photograph a sampling weight of one divided by
the number of selected photographs from its city. The weights within every
city sum to one, giving every city a sampling probability of 20\%.
The official weighted sampler draws with replacement. The full original
selection of 44,257 paths and its order were retained.

The following counts were calculated by the official sampler using seed 0
and the full budget. The total in each column is 249,984 presentations.
They show how the experiment redistributed a fixed amount of training.

\begin{center}\small
\begin{tabular}{lrr}
\toprule
City & Starter sampling & Balanced sampling \\
\midrule
Seoul & 63,646 & 50,153 \\
Busan & 29,328 & 49,745 \\
Tokyo & 75,426 & 50,023 \\
Osaka & 40,670 & 49,939 \\
Taipei & 40,914 & 50,124 \\
\bottomrule
\end{tabular}
\end{center}

Every selected photograph appeared at least once with starter sampling.
With weighted sampling and replacement, 43,818 distinct photographs appeared
at least once. These counts describe the actual draws from the full selection.

\subsection{Local execution and controlled comparison}

The first MPS smoke check failed during backward computation with a tensor
stride/view error. Converting the MPS input batch to contiguous memory in
\texttt{recipe.mix} allowed the official training, validation, and ONNX
export to complete. This operation preserves the input values and labels.
Both local conditions used this same compatibility adjustment.
The official training, rules, export, and submission-check source files
and the pinned package versions were preserved.

Both conditions used the same MPS device, seed 0, worker count 0, 224
$\times$ 224 inputs, batch size 64, 3,906 steps, and 249,984 presentations.
Cropping, flipping, color jitter, AdamW, OneCycleLR, and label smoothing
were the same. The official MPS path used float32. The earlier CUDA runs
used mixed precision, which is one reason to measure a fresh local control.

The control recipe was committed as \texttt{a108f5c}, with recipe hash
\texttt{7c1cb9f0ee41}. The balanced recipe was committed as \texttt{bb6bc15},
with recipe hash \texttt{e7500546f0a5}. Both recorded the full-selection
hash \texttt{1a1ca6c18823}. Saved snapshots matched their recorded commits,
and each result was checked against its experiment CSV row and confusion
matrix. Both conditions and the exact final recipe passed official smoke
checks. Smoke checks are recorded separately from full-budget training.

\begin{center}\small
\begin{tabular}{llrr}
\toprule
Run tag & Seed & Training time (min) & Validation (\%) \\
\midrule
\texttt{baseline\_mps} & 0 & 47.2 & 49.51 \\
\texttt{balanced} & 0 & 48.0 & 50.91 \\
\bottomrule
\end{tabular}
\end{center}

Training took 2,829.9 and 2,881.0 seconds, respectively, with validation and
export taking additional time. Both runs satisfied the official budget.

\subsection{Results and the student\textquotesingle{}s interpretation}

\begin{center}\small
\begin{tabular}{lrrr}
\toprule
City & Control (\%) & Balanced (\%) & Difference (pp) \\
\midrule
Seoul & 57.41 & 51.33 & -6.08 \\
Busan & 39.57 & 46.62 & +7.05 \\
Tokyo & 61.63 & 50.01 & -11.62 \\
Osaka & 35.81 & 44.97 & +9.16 \\
Taipei & 53.14 & 61.61 & +8.47 \\
\midrule
Mean over cities & 49.51 & 50.91 & +1.40 \\
\bottomrule
\end{tabular}
\end{center}

Busan, Osaka, and Taipei increased, while Seoul and Tokyo decreased.
The increases outweighed the decreases in the equally weighted city mean.
After the agent explained the method and the score calculation, the student
interpreted the Busan increase as more frequent exposure giving the model
more opportunities to learn characteristics of Busan. The observed increase
from 39.57\% to 46.62\% is consistent with that interpretation. Identifying
the particular visual characteristics used requires further evidence.

This experiment shows that changing how frequently the existing photographs
are sampled can affect the measured score, with gains and losses differing
by city. Each local condition was run once with seed 0. The improvement
of 1.40 percentage points describes this matched pair, and repeated-seed
variability for balanced sampling remains unmeasured. The original CUDA
baseline mean of 50.14\% remains a separate two-seed result.

The agent selected balanced sampling for final Track A release
\texttt{track-a-2}, based on its higher score in the matched local comparison.
The recipe and full selection were released together. The appendix contains
the sampling audit, original run logs, and comparison evidence.

\section{Track B experiments and submission}

\subsection{First model training and release}

After reviewing the erasing results, the student selected preparation of the
first Track B submission as the next step. The experiment \texttt{b\_baseline}
used the provided \texttt{src/train\_b.py}, with ResNet18 initialized from
ImageNet-1k V1, 224 $\times$ 224 inputs, all starter training photographs,
and four epochs. The default configuration used batch size 64, AdamW with
learning rate 0.0003 and weight decay 0.05, OneCycleLR, and label smoothing 0.1.

This script trains independently of the Track A recipe and evaluates the
full starter validation set. The first run used the provided script as
supplied. It had no seed-setting argument, so the experiment record has an
empty seed field. This difference is relevant when comparing later runs.
The run presented 176,896 photographs and recorded 1,637.1 seconds of
training, approximately 27.3 minutes. It ran on CUDA from commit
\texttt{50530fc} and finished at 01:14 KST on 9 November.

\begin{center}\small
\begin{tabular}{lr}
\toprule
City & Validation accuracy (\%) \\
\midrule
Seoul & 55.40 \\
Busan & 37.55 \\
Tokyo & 63.77 \\
Osaka & 34.87 \\
Taipei & 53.14 \\
\midrule
Mean across cities & 48.95 \\
\bottomrule
\end{tabular}
\end{center}

The \texttt{--submit} option generated \path{submission/model.onnx} and
\path{submission/submission.json}. The model passed the official format and
Track B size-limit checks on Colab and locally. Local measurements were
11,174,283 parameters, 1.813563904 GMACs, and 44.707595 MB, with ONNX opset 17.
CPU inference at batch sizes 1 and 3 returned five finite logits for each
photograph. The validation score above is the value recorded during Colab
training.

The release \texttt{track-b-1} was published at 01:22:42 KST on 9 November.
Its title was \texttt{Starter ResNet18 at 224, 4 epochs}. Release publication,
both completed asset uploads, and the agreement between the GitHub asset
SHA-256 values and the checked files were verified. The release commit was
\texttt{d4f7ce3}, and the training commit was \texttt{50530fc}. Full hashes and
asset sizes are retained in the appendix evidence and release-query JSON.

\subsection{Scope and preparation of the student\textquotesingle{}s additional photographs}

The student said that Seoul and Busan photographs were available in Downloads.
When the agent initially proposed processing all IMG photographs, the student
restricted the scope to files downloaded that day. Finder\textquotesingle{}s date-added
metadata, converted to Korean time, identified seven files added at
01:27:30 KST on 9 November. The student confirmed that
\texttt{IMG\_0549.HEIC} was Seoul and the other six photographs were Busan.

\begin{center}\small
\begin{tabular}{llr}
\toprule
City & Original photograph numbers & Count \\
\midrule
Seoul & 0549 & 1 \\
Busan & 0237, 0240, 0276, 0286, 1403, 1413 & 6 \\
\bottomrule
\end{tabular}
\end{center}

The HEIC images were converted into orientation-corrected RGB copies with
512 pixels on the shorter side. The provided photo-preparation code saved
city-organized JPEGs, and checks confirmed empty EXIF metadata. All seven
files had distinct file hashes and decoded-pixel hashes. The SHA-256 of each
copy was recorded in a manifest. The collection scope was the student\textquotesingle{}s one
Seoul photograph and six Busan photographs supplied on 9 November.

Because the first Track B run had an unspecified seed, the agent added an
optional \texttt{--seed} argument for the data comparison. Both new conditions
used seed 0, ImageNet V1 ResNet18, 224 $\times$ 224 inputs, four epochs, and
batch size 64. They used 44,257 and 44,264 training photographs, respectively.
Both yielded 691 complete batches per epoch and 176,896 photograph
presentations per run.

A photo-inventory check verified file paths and hashes before training.
Synthetic-data checks covered repeatability with seed 0, experiment recording,
and compatibility with the original unspecified-seed behavior. Both full
runs also passed the exact seven-file inventory and hash checks before
training began.

\subsection{Paired local comparison on the Mac}

Colab\textquotesingle{}s GPU usage limit prevented a new GPU connection, so the work moved to
the student\textquotesingle{}s Mac at the student\textquotesingle{}s request. The environment was an Apple
M1 Pro with 16 GiB memory and macOS 15.6, running PyTorch 2.6.0 and torchvision
0.21.0 with MPS operations verified. The 44,257 starter training photographs
and 12,077 validation photographs were copied into the local working folder.
The training paths and order matched the existing all-photo selection.

Both conditions ran the same \texttt{src/train\_b.py} with seed 0 and zero
loader workers. They used ImageNet V1 ResNet18, 224 $\times$ 224 inputs, four
epochs, batch size 64, AdamW, OneCycleLR, and the original training
augmentations. Both full runs used commit \texttt{500ed72}.

An initial short check with two workers was delayed during process shutdown
after saving its outputs. With zero workers, training, export, checking, and
normal shutdown were verified before full training began. These short checks
were recorded separately from the eleven full experiments.

The control tag was \path{b_control_seed0}, and the added-photo tag was
\path{b_own_sb_seed0}. Seoul training photographs increased from 11,274 to
11,275, and Busan photographs from 5,191 to 5,197. The other three cities
and the validation set had the same inventories in both conditions.

\begin{center}\small
\begin{tabular}{p{0.30\linewidth}rrr}
\toprule
Condition & \shortstack{Training\\photographs} &
\shortstack{Photograph\\presentations} & \shortstack{Validation\\(\%)} \\
\midrule
Starter photographs & 44,257 & 176,896 & 48.61 \\
Add 7 photos & 44,264 & 176,896 & 48.74 \\
\bottomrule
\end{tabular}
\end{center}

The recorded training times were 2,091.5 seconds, approximately 34.9 minutes,
for the control and 2,098.4 seconds, approximately 35.0 minutes, for the
added-photo condition. Validation and ONNX export took additional time.
The runs finished at 03:16:31 and 03:52:58 KST on 9 November, respectively.

\subsection{Added-photo results and interpretation}

\begin{center}\small
\begin{tabular}{lrrr}
\toprule
City & Starter photos (\%) & Add seven (\%) & Change (pp) \\
\midrule
Seoul & 61.44 & 56.34 & -5.10 \\
Busan & 36.91 & 37.19 & +0.28 \\
Tokyo & 59.86 & 61.63 & +1.77 \\
Osaka & 34.02 & 35.43 & +1.41 \\
Taipei & 50.82 & 53.10 & +2.28 \\
\midrule
Mean across cities & 48.61 & 48.74 & +0.13 \\
\bottomrule
\end{tabular}
\end{center}

These figures use the scores stored to four decimal places in
\texttt{experiments.csv}. Among the cities receiving additional photographs,
Seoul accuracy decreased and Busan accuracy increased slightly. Tokyo, Osaka,
and Taipei accuracies also changed. The overall observed increase was
0.13 percentage points. One run per condition at one seed provides limited
information about persistent improvement and training variability.

Adding photographs changes the dataset length and shuffled order. Connecting
each per-city change directly to an individual added photograph requires
further verification. The appropriate control for this paired data experiment
is the 48.61\% result from the same device and seed. A comparison with the
first Colab model\textquotesingle{}s 48.95\% also includes device and seed-setting differences.

The seven photographs represent approximately 0.016\% of the starter training
set. Generalizing from this small addition to the effect of collecting many
more photographs requires experiments with greater numbers and diversity.
The numerical comparisons and limitations in this section are the agent\textquotesingle{}s
analysis of the records. Additional student interpretations can be recorded
with their authorship identified.

Both ONNX models passed the official Track B format and limit checks. Each
had 11,174,283 parameters, 1.813563904 GMACs, a file size of 44,707,595 bytes,
and opset 17. Full ONNX checking and CPU inference at batch sizes 1 and 3
also confirmed five finite outputs per photograph.

Each model and its metadata were stored in a separate result folder. The
queued \texttt{--submit} command initially placed the seven-photo model in
the local \texttt{submission/} folder after training. The submission review
on 9 November selected the first model and restored it to the submission
folder. The selection was revised on 10 November as described below.
Model hashes, original run outputs, and check results are included in the
appendix evidence.

\subsection{Using the balanced Track A model for Track B}

The leaderboard snapshot updated at 17:04 KST on 10 November reported
38.26\% for \texttt{track-b-1}, rank 26 among 27 scored models. After
expressing concern about that result, the student selected the agent\textquotesingle{}s
proposal to validate and submit the already trained balanced model.

The source was \path{results/balanced_seed0/model.onnx}, trained from
commit \texttt{bb6bc15} on MPS with seed 0, batch size 64, 3,906 steps,
and 249,984 presentations. We evaluated all 12,077 starter validation
photographs with the official Track B checker using ONNX Runtime on CPU.
The exact mean city accuracy was 0.509104010051048, or 50.91\%.
Its confusion matrix matched the original balanced MPS validation result.

\begin{center}\small
\begin{tabular}{lrr}
\toprule
City & First B model, ONNX/CPU (\%) & Balanced, ONNX/CPU (\%) \\
\midrule
Seoul & 55.40 & 51.33 \\
Busan & 37.55 & 46.62 \\
Tokyo & 63.75 & 50.01 \\
Osaka & 34.87 & 44.97 \\
Taipei & 53.14 & 61.61 \\
\midrule
Mean across cities & 48.94 & 50.91 \\
\bottomrule
\end{tabular}
\end{center}

The increase of 1.97 percentage points supports selecting the balanced
model from these available candidates. The two models also differ in
training budget, device, and seed-setting conditions. The paired sampling
comparison in Section 7 provides the controlled evidence for the sampling
change, with 49.51\% and 50.91\% under the same MPS and seed-0 conditions.
Reusing a completed model added no training run to the eleven recorded
full experiments.

The checked model was published as \texttt{track-b-2} at 17:32:25 KST
on 10 November, titled \texttt{ResNet18 at 224, balanced sampling, seed 0}.
Both release assets were downloaded without authentication and matched
the validated local files byte for byte. The model has 11,174,283 parameters,
1.813563904 GMACs, and a size of 44,707,595 bytes. The complete model and
metadata hashes and the official full check are reproduced in the appendix.

The subsequent leaderboard update at 17:45 KST reported 42.03\% for this
release, rank 14 among 27 scored models. Compared with the first release\textquotesingle{}s
38.26\%, the public score increased by 3.77 percentage points. It also
exceeded the displayed starter-model reference of 40.3\% by approximately
1.7 percentage points. The leaderboard and local validation use different
photographs, so their absolute scores describe separate evaluation sets.

\begin{center}\small
\begin{tabular}{lrr}
\toprule
City & First B leaderboard (\%) & Balanced B leaderboard (\%) \\
\midrule
Seoul & 58.0 & 62.3 \\
Busan & 15.9 & 30.9 \\
Tokyo & 70.5 & 61.8 \\
Osaka & 23.7 & 30.0 \\
Taipei & 23.2 & 25.1 \\
\midrule
Published overall score & 38.26 & 42.03 \\
\bottomrule
\end{tabular}
\end{center}

The per-city public values are rounded in the source data. The overall
row uses each published score, retaining its provided precision.

\section{Post-training occlusion analysis and interpretation checks}

\subsection{Case selection and measurement}

The student requested an optional analysis of which image regions affected
the model\textquotesingle{}s outputs. This analysis was completed on 9 November with
the first release, \texttt{track-b-1}, which was selected at that time.
The later balanced-model submission uses a different ONNX file.
The agent shuffled each city\textquotesingle{}s validation list with seed
0 and evaluated the first 32 photographs per city, 160 in total. In that
order, the first correctly classified and first incorrectly classified
photograph from each city were selected, producing ten cases. This sampling
supports a comparison of correct and incorrect cases. Overall accuracy is
reported from the full validation set.

The official \texttt{rules.grader\_input} produced a 224 $\times$ 224 center
crop. Each input was partitioned into 49 non-overlapping 32 $\times$ 32
regions. One region at a time was filled with RGB (0.5, 0.5, 0.5), followed
by inference. For each photograph, the original predicted city was fixed as
the target, and we measured the change in that target\textquotesingle{}s logit after masking.
A logit is the output score before conversion to a softmax probability.

Red indicates a decrease in the original predicted city\textquotesingle{}s logit after
masking, and blue indicates an increase. All ten heatmaps use the same color
scale. The probability columns below also refer to the original predicted
city.

\begin{center}\small
\begin{tabular}{lllrrr}
\toprule
City & Original case & \shortstack{Original\\prediction} &
\shortstack{Original\\prob. (\%)} & \shortstack{Masked\\prob. (\%)} &
\shortstack{Prediction\\changes /49} \\
\midrule
Seoul & Correct & Seoul & 65.0 & 42.9 & 1 \\
Seoul & Incorrect & Busan & 55.1 & 41.1 & 0 \\
Busan & Correct & Busan & 40.6 & 31.1 & 1 \\
Busan & Incorrect & Tokyo & 34.9 & 29.0 & 9 \\
Tokyo & Correct & Tokyo & 85.1 & 78.4 & 0 \\
Tokyo & Incorrect & Seoul & 64.8 & 48.7 & 0 \\
Osaka & Correct & Osaka & 93.0 & 90.6 & 0 \\
Osaka & Incorrect & Tokyo & 45.4 & 32.2 & 4 \\
Taipei & Correct & Taipei & 39.0 & 28.4 & 0 \\
Taipei & Incorrect & Seoul & 55.5 & 38.1 & 0 \\
\bottomrule
\end{tabular}
\end{center}

The masked probability is measured at the region giving the largest decrease
in the original predicted city\textquotesingle{}s logit. The final column counts how many of
the 49 masks changed the predicted city. Maximum logit decrease and a change
in the final predicted city were recorded as separate measures.

Across 490 masked inputs, 15 predictions changed, occurring in four of the
ten photographs. One mask changed each of the correctly classified Seoul
and Busan cases to an incorrect prediction. Four masks changed the
incorrect Osaka case to the correct city, Osaka.

The numerical results, original logits, case list, analysis code, and figures
were retained. Checks covered the fixed original target, complete input
coverage by the 49 regions, model and photograph hashes, and agreement
between the stored arrays and CSV files.

\subsection{The student\textquotesingle{}s question and replacement controls}

In the incorrectly classified Osaka photograph, masking a sky region changed
the prediction to Osaka. The student questioned this result in the following
statement, translated from Korean.

\begin{quote}\small
Masking a completely unexpected area changes the prediction, so I really do
not understand why.
\end{quote}

In response, we compared different replacement methods in the same photograph
at the same coordinates. The region was \texttt{x [64,96)} and
\texttt{y [96,128)} in the model input crop. The conditions included the mean
color of the original sky patch, Gaussian blur with radius 8 pixels applied
only to that patch, and grayscale fills at different brightness levels.

\begin{center}\small
\begin{tabular}{lrrl}
\toprule
Replacement & Tokyo prob. (\%) & Osaka prob. (\%) & Prediction \\
\midrule
Original & 45.4 & 19.9 & Tokyo \\
Middle gray 0.50 & 32.2 & 34.7 & Osaka \\
Original sky patch mean & 43.8 & 21.2 & Tokyo \\
Blur, radius 8 pixels & 45.6 & 19.9 & Tokyo \\
Dark gray 0.25 & 36.1 & 29.9 & Tokyo \\
Light gray 0.75 & 32.5 & 36.8 & Osaka \\
\bottomrule
\end{tabular}
\end{center}

The predicted city depended on how the same sky region was replaced. The
initial gray-mask result combined a change to the original information with
the introduction of a new patch color and boundary. This case demonstrates
output sensitivity to particular masking conditions. Interpreting semantic
or cultural city cues requires considering this dependence on the analysis
method. We also consulted Section 4.3 of Fong and Vedaldi\textquotesingle{}s ICCV 2017 paper,
which discusses artifacts introduced by sharp color contrasts from constant
perturbations.

This diagnostic changed the inputs to a completed model. Its case-level
outputs and replacement controls are stored separately from the eleven full
training experiments. The report distinguishes training changes, diagnostic
input changes, observed output sensitivity, and hypotheses about semantic
causes.

\section{Conclusions from the measurements and further work}

Among the three initial CUDA Track A conditions, the highest mean validation score was
50.14\%, obtained with the baseline recipe and all training photographs.
The student\textquotesingle{}s 56-photo exclusion scored 49.26\%, and the student-proposed
erasing experiment scored 49.30\%. Both changes scored lower than the
baseline with the same two seeds. These conclusions apply to the particular
selection and augmentation settings tested.

The later matched MPS comparison improved from 49.51\% with starter sampling
to 50.91\% with balanced sampling. Redistributing training appearances
increased accuracy for Busan, Osaka, and Taipei and reduced it for Seoul and
Tokyo. The student interpreted more frequent exposure as giving the model
more opportunities to learn Busan characteristics. The measured score
improvement applies to the one completed seed-0 pair. Balanced sampling
was selected for the final Track A release.

Photographs that seem difficult for a person to identify can still affect
training. The intended benefit of encouraging multiple cues through
augmentation also required measurement. Further analysis is needed to
identify the specific causes of the decreases and the model\textquotesingle{}s use of visual
features. The post-training diagnostic demonstrated sensitivity both to
which region was masked and to how it was replaced, illustrating the
importance of the analysis conditions when interpreting a model.

In Track B, the student\textquotesingle{}s decision to supply one Seoul photograph and six
Busan photographs was evaluated against a control on the same MPS device
and seed. The observed change was +0.13 percentage points. Repeated
measurements are needed to interpret this small difference from a single
pair of runs. On 10 November, the student selected reuse of the balanced
Track A model for Track B. Full ONNX/CPU validation increased from 48.94\%
for the first B model to 50.91\%. The subsequent public leaderboard score
increased from 38.26\% to 42.03\%. Differences in training budget, device,
and seed-setting conditions also contribute to this model-selection
comparison. The controlled evidence for balanced sampling comes from the
matched Track A pair.

With one more week, the agent would propose additional seeds under matching
conditions and a larger, more varied collection of street, residential,
and transport scenes from each city. Gradually increasing the collection
would allow a comparison of the effect of data quantity. The 56-photo
exclusion could also be divided into more specific reason-based groups to
investigate which groups relate to score changes. Specific choices of
additional collection and experiments would be made by the student and
recorded with their rationale.

\section{Final submission and validation checks}

The selected releases use the completed balanced ResNet18 model for both
tracks. The eleven completed full training runs and their measurement
limitations are reported above. MobileNet experiments remain in progress.

\subsection{Validation and leaderboard scores}

\begin{center}\small
\begin{tabular}{llrl}
\toprule
Track & Final release & Validation (\%) & Leaderboard \\
\midrule
A & \texttt{track-a-2} & 50.91, MPS seed 0 & Pending \\
B & \texttt{track-b-2} & 50.91, ONNX/CPU & 42.03\% \\
\bottomrule
\end{tabular}
\end{center}

The public data checked at 17:54:05 KST on 10 November showed an update
time of 17:45 KST. It contained \texttt{SoSleepy} with \texttt{track-b-2}
and rank 14 of 27. The Track A list still contained no matching nickname,
so its official leaderboard score remains unavailable. The original
snapshot and extracted status are retained in
\path{outputs/leaderboard_status.json} and
\path{outputs/leaderboard_1745.json}.

The leaderboard uses YouTube street-video frames from sources separate from
starter validation. Its Track A result averages seeds 0, 1, and 2, while
our final balanced result uses MPS seed 0. Source, scene, and execution
differences can affect the scores. The local balanced result of 50.91\%
and the B leaderboard result of 42.03\% therefore refer to different
evaluation sets. The official A result still requires a public grading entry.

\subsection{Submission file checks}

The final Track A recipe passed official rules and the 20-step MPS
\texttt{--quick} check, including validation and export. The 44,257-photo
selection retained all original paths and their order. The downloaded
release files matched the checked local files.

The selected balanced Track B ONNX file passed official format and
size-limit checks and full CPU validation on all 12,077 photographs.
Its exact score was 0.509104010051048, or 50.91\%. Each checked local
submission file matched the corresponding published \texttt{track-b-2}
asset SHA-256. The appendix includes the evidence.

\section*{References}
\begin{enumerate}\footnotesize
\item \href{https://github.com/jdasam/mas1004-assignment2}{Assignment README and official training code, GitHub}.
\item \href{https://github.com/jellykeeng/mas1004-assignment2/releases/tag/track-a-2}{Student repository, final Track A release}.
\item \href{https://github.com/jellykeeng/mas1004-assignment2/releases/tag/track-b-2}{Student repository, selected balanced Track B release}.
\item \href{https://github.com/jellykeeng/mas1004-assignment2/releases/tag/track-b-1}{Student repository, first Track B release used for the occlusion analysis}.
\item \href{https://jdasam.github.io/mas1004-leaderboard/board.json}{Public leaderboard data}.
\item \href{https://news.seoul.go.kr/safe/archives/515090}{Seoul Metropolitan Government, overhead-line undergrounding plan update, 28 October 2024}.
\item \href{https://docs.pytorch.org/vision/0.21/generated/torchvision.transforms.RandomErasing.html}{torchvision 0.21 documentation for RandomErasing}.
\item \href{https://captum.ai/api/occlusion.html}{Captum documentation for Occlusion}.
\item \href{https://robots.ox.ac.uk/~vedaldi/assets/pubs/fong17interpretable.pdf}{Fong and Vedaldi, Interpretable Explanations of Black Boxes by Meaningful Perturbation, ICCV 2017}.
\end{enumerate}

\clearpage
\appendix
\section{Photo review material}

\begin{figure}[htbp]
\centering
\picturefile{figures/mistakes_baseline_seed0_busan_00.png}
\caption{Forty-eight high-confidence mistakes among photographs labeled Busan
under the seed 0 model. Predictions and output probabilities appear below
the photographs. The student specifically identified roadside parking in
photograph 8 and utility poles and wires in photograph 48.}
\end{figure}

\clearpage
\begin{figure}[htbp]
\centering
\picturefile{figures/mistakes_baseline_seed0_osaka_00.png}
\caption{Forty-eight high-confidence mistakes among photographs labeled Osaka
under the seed 0 model. The student\textquotesingle{}s initial observations concerned shrines,
houses, and residential streets with no visible sidewalks.}
\end{figure}

\clearpage
\section{Occlusion analysis figures}

\begin{figure}[!ht]
\centering
\picturefile{figures/occlusion_report_examples.png}
\caption{A correctly classified Seoul case and an incorrectly classified
Osaka case from the first Track B model, \texttt{track-b-1}.
The original predicted city is fixed as the target. Black
outlines indicate the regions with the largest target-logit decrease.
Red indicates a logit decrease after masking and blue an increase.
All ten cases use the same color scale.}
\end{figure}

\clearpage
\begin{figure}[!ht]
\centering
\picturefile{figures/occlusion_report_controls.png}
\caption{Replacement controls for the first Track B model on the same
sky region of the incorrectly classified Osaka photograph. Middle and light gray produced Osaka
predictions. The original sky patch mean and blur produced Tokyo predictions.
This case demonstrates output sensitivity to the replacement method.}
\end{figure}

\clearpage
\section{Recorded outputs and numerical evidence}

Original machine outputs are reproduced in their recorded form. The evidence
files support the numerical tables in the main report.

\subsection*{All full training runs}
\outputfile{outputs/experiments.csv}

\subsection*{Original format checks for the six Track A models}
\outputfile{outputs/1105_155816_check_submission.txt}
\outputfile{outputs/1105_165326_check_submission.txt}
\outputfile{outputs/1108_215154_check_submission.txt}
\outputfile{outputs/1108_223906_check_submission.txt}
\outputfile{outputs/1108_235701_check_submission.txt}
\outputfile{outputs/1109_003754_check_submission.txt}

\subsection*{Per-city summary lines from mistake-sheet generation}
\outputfile{outputs/curate_summary.txt}

\clearpage
\subsection*{Values calculated from the original baseline results}
The supplied \texttt{make\_evidence.py} calculated these values from result
JSON files, sheet CSV files, and the baseline recipe. Input files and full
run logs are retained in the project\textquotesingle{}s \texttt{outputs/} folder.
\outputfile{outputs/analysis_summary.txt}

\clearpage
\subsection*{Evidence for the exclusions and erasing configuration}
The evidence script used four full training results, the selection lists,
the candidate CSV files, and the new recipe.

Selected training paths are recorded in the following file.
\par\path{outputs/curated56_selection.csv}\par
The exclusion plan is recorded in the following file.
\par\path{outputs/curated56_curation_plan.csv}\par
\outputfile{outputs/curation_analysis.txt}

\clearpage
\subsection*{Erasing results and comparison of the three conditions}
These values were calculated from the six original Track A results, run
records, recipes, and full-photo selection.
\outputfile{outputs/erasing_analysis.txt}

\clearpage
\subsection*{Original first Track B training and format-check outputs}
\outputfile{outputs/1109_004557_train_b_b_baseline.txt}
\outputfile{outputs/1109_011454_check_submission.txt}

\subsection*{Track B records and public release verification}
This summary was calculated from the Colab experiment row, submission
metadata, local measurements, and public release and tag query results.
\outputfile{outputs/track_b_analysis.txt}

\clearpage
\subsection*{Student-photo preparation and paired-comparison evidence}
This summary used the date metadata of the files supplied on 9 November,
the city labels confirmed by the student, the actual prepared inventory,
and local verification results. Both paired full training runs were completed.
\outputfile{outputs/personal_photo_analysis.txt}

\clearpage
\subsection*{Original local MPS training and format-check output}
\outputfile{outputs/local_training_mps_20261109.log}

\subsection*{MPS scores and verification evidence}
This summary cross-checked the two experiment CSV rows, comparison result
JSON, both submission metadata files, actual run log, and photo-inventory
checks. The original JSON is retained in
\path{outputs/local_mps_comparison_seed0.json}.
\outputfile{outputs/local_mps_analysis.txt}

\clearpage
\subsection*{Submission checks and selection recorded on 9 November}
These checks preceded the balanced-sampling release. The earlier MPS
smoke-check error is retained in
\path{outputs/final_track_a_quick_mps_failed.txt}.
\outputfile{outputs/final_track_a_quick.txt}
\outputfile{outputs/final_track_b_check.txt}
\outputfile{outputs/final_selection_analysis.txt}

\clearpage
\subsection*{Balanced sampling, matched local comparison, and final release}
\outputfile{outputs/balanced_comparison.txt}
\outputfile{outputs/balanced_sampling_audit.txt}
\outputfile{outputs/baseline_mps_training.txt}
\outputfile{outputs/balanced_training.txt}
\outputfile{outputs/balanced_final_quick.txt}
\outputfile{outputs/balanced_release_verification.txt}

\clearpage
\subsection*{Balanced-model reuse and the second Track B release}
The official checker evaluated all 12,077 validation photographs. The public
release files were retrieved anonymously and matched the checked candidates.
\outputfile{outputs/balanced_b_full_check.txt}
\outputfile{outputs/balanced_b_selection.txt}
\outputfile{outputs/balanced_b_release_verification.txt}
\outputfile{outputs/leaderboard_1745_analysis.txt}

\clearpage
\subsection*{Occlusion measurements and replacement-control evidence}
This summary cross-checked the analysis JSON, case-level and region-level
CSV files, and original logit arrays. The original arrays are retained in
\path{outputs/occlusion_raw_outputs.npz}.
\outputfile{outputs/occlusion_analysis.txt}

\clearpage
\section{Lab notebook in its original language}

The complete dated entries of \texttt{LOG.md} are reproduced below in their
original language, including the student\textquotesingle{}s Korean statements. English
translations of the relevant observations appear in the main report.
The original file including template guidance is retained in
\path{outputs/LOG_original.md}.
\outputfile{outputs/LOG.md}

\end{document}
